\documentclass[12pt]{article}
\usepackage[utf8]{inputenc}
\usepackage[T1]{fontenc}
\usepackage{hyperref}
\title{Taushiro Language Resources for K-12 Curriculum Development}
\author{Compiled by Claude Sonnet 5 (claude-sonnet-5) for Andr\'as Kornai}
\date{}
\begin{document}
\maketitle

\section{Introduction}
Taushiro (ISO 639-3 trr; Glottocode taus1253) is a critically endangered language isolate of Peru. Sources found in search indicate that only one speaker, Amadeo Garc\'ia Garc\'ia, remained. Very little published material exists, and no school materials were found. A K-12 curriculum in the language is not currently feasible from existing resources.

\section{Primary Educational Resources}
No primers, textbooks or school programmes were found in search results. A Peruvian news item reports the official recognition of alphabets for three Amazonian languages; check it for whether Taushiro is among them:
\url{https://www.actualidadambiental.pe/reconocen-oficialmente-alfabetos-de-tres-lenguas-originarias-de-la-amazonia/}

\section{Linguistic and Reference Materials}
\begin{itemize}
\item Berkeley SAPHON phonological inventory of Taushiro (consonant and vowel chart; cites Alicea 1975 and O'Hagan's fieldnotes): \url{https://linguistics.berkeley.edu/saphon/en/inv/Taushiro.html}
\item Max Planck EVA page on Taushiro numerals: \url{https://lingweb.eva.mpg.de/channumerals/Taushiro.htm}
\item Endangered Languages Project bibliography for Taushiro: \url{https://endangeredlanguages.com/lang/2163/bibliography?hl=pt-br}
\item Wikipedia overview of the language: \url{https://en.wikipedia.org/wiki/Taushiro_language}
\end{itemize}
Search summaries mention Alicea's \emph{Vocabulario taushiro} (SIL, 1975) as the main wordlist, but no online copy was seen.

\section{Religious and Community Resources}
None found.

\section{Digital and Audio Resources}
Search summaries mention a Ministry of Culture database of about 1,500 words, 27 stories and three songs recorded with the last speaker. No public URL for it was seen. Background on the speaker: \url{https://en.wikipedia.org/wiki/Amadeo_Garc%C3%ADa_Garc%C3%ADa}

\section{Curriculum Development Context}
With one reported speaker, a curriculum would have to be revitalization or documentation oriented and based on the archived recordings. Dates and names in the secondary sources conflict, so verify them against the originals.
\end{document}